\section{本章小结}

本章把微积分中“沿长度、体积累计”的想法推广到任意对象组成的空间，得到五条相互衔接的认识。

第一，测量分两层，而且两层互不相同。$\sigma$代数规定\emph{哪些集合看得见}，测度规定\emph{看得见的集合有多大}；还要把测度与度量分清：测度吃进一个集合给出份量，度量吃进两个点给出距离，同一底集上两者各自独立选择，讨论Wasserstein等几何问题时才在测度之外另加度量。同一个可测范围上可以配任意多种权重，也可以通过子$\sigma$代数在不删点的情况下降低分辨能力。有限、概率、$\sigma$有限是三种强度不同的质量条件，零测集和完备化说明哪些点上的例外可以忽略。当开集参与进来时，Borel结构、Polish空间和支持把“测量”与“位置”联系起来，而一般积分、三条收敛定理、乘积测度与Tonelli/Fubini都在明确条件下原样沿用。

第二，同一空间上的测度比较从零集关系开始，再到表示和定量指标。绝对连续、等价和相互奇异刻画零集之间的包含、重合与分离；Radon–Nikodym密度是两种称重方式的局部换算比例，Lebesgue分解把能这样换算的部分和集中在零集上的部分分开。为描述测度之差引入了有符号测度及其Jordan分解与总变差。全变差、Hellinger和KL只比较同一位置上的事件或密度，Wasserstein还要计入两点间的距离；积分概率度量把它们统一成“某类探针上的最大读数差”。这些指标回答的问题不同，不存在一个能代替其余的“万能距离”。

第三，跨空间时质量的传递规则从确定走向随机。可测映射的推前回答“每个输出事件收到多少输入质量”，输出积分等于输入上的复合积分；测度核把“一个输入对应一个点”放宽为“一个输入对应一项目标测度”，核积分和核复合分别统一了输出传递与观测传递。联合测度保留输入输出的配对，乘积和确定图是其特例；在标准Borel空间上，联合测度可以沿第一边缘分解为归一化条件核。条件积分只保存粗信息下的累计，条件期望等概率法则在后文展开。

第四，跨空间比较前必须先固定对应关系。给定映射或核以后，比较的是变换后的输出而不是原对象；表示可能丢失信息，所以输出相同不能反推输入相同，不同表示之间也没有收缩性质。给定跨空间成本时，Kantorovich计划允许拆分质量、Monge映射不允许，二者可行域不同；可行、有限、最优值达到要分别靠耦合非空、有限成本候选和紧性加下半连续性来保证。输出匹配为零与运输代价非零可以同时成立。

第五，测度序列是否趋近取决于探针。全变差统一控制所有事件，弱收敛只看有界连续函数，Wasserstein还要加上距离矩的尾部控制；弱收敛通过Portmanteau界读取边界零质量事件，在连续映射下保持，在下半连续成本下只给积分下界。一致紧性经Prokhorov定理保证弱子列，并支撑最优运输计划的存在性；固定时刻、有限时域一致和长时间极限是三种不能互换的逼近。
这些通用操作为后文的随机过程、条件模型、经验测度和策略分析提供了语言，但相应的概率性质、统计保证和生成机制仍需各自的条件。
